\chapter{Implementation}
\label{chap:implementation}

\section{Overview}
\label{sec:impl-overview}

This chapter describes the concrete implementation of the framework laid out in
Chapter~\ref{chap:system-design}. The codebase is hosted on GitHub at
\url{https://github.com/Vivian-Lopez/Monte-Carlo-and-Automatic-Differentiation-Systematic-Benchmarking-Framework}
and is organised as a single Python package with optional native subpackages.
The top-level directory structure is:

\begin{verbatim}
benchmarking/
    workloads/        # WorkloadConfig dataclasses + registry
    engines/          # Python wrappers for each engine
    cpp/              # native C++ implementation (pybind11)
    rust/             # native Rust implementation (PyO3)
    runner.py         # BenchmarkRunner
    db.py             # BenchmarkDB (SQLite)
    api/              # Flask REST API
frontend/             # React + TypeScript + Vite
experiments/          # CLI experiment drivers
results/              # SQLite database, generated tables
tests/                # pytest suite
\end{verbatim}

The Python package is pure-Python apart from optional native extensions. The C++
extension is built with \texttt{pip install -e benchmarking/cpp/} (using a
\texttt{setup.py} that invokes pybind11); the Rust extension is built with
\texttt{maturin develop --release} inside \texttt{benchmarking/rust/}. If either
extension is not built, the corresponding engine is silently excluded from the
engine registry at runtime, so the framework remains fully functional on systems
without a C++ or Rust toolchain.

\section{Core Modules}
\label{sec:core-modules}

\subsection{Workload Registry}
\label{sec:impl-registry}

The workload registry is implemented in \texttt{benchmarking/workloads/base.py}.
The decorator \texttt{@workload(name)} performs three actions in a single pass:
it applies \texttt{dataclass} to the decorated class, sets a class-level
attribute \texttt{workload\_type = name}, and registers the resulting class in
the module-level dictionary \texttt{WORKLOAD\_REGISTRY}. The base class
\texttt{WorkloadConfig} provides concrete implementations of \texttt{validate()},
\texttt{to\_dict()}, \texttt{from\_dict()} and \texttt{config\_hash()} so that no
subclass is required to override anything except its field list and (optionally)
its validation logic.

The free function \texttt{config\_from\_dict(d)} reads \texttt{d["workload\_type"]},
looks up the corresponding class in \texttt{WORKLOAD\_REGISTRY}, and delegates
to that class's \texttt{from\_dict()}. Unknown keys are silently dropped during
deserialisation, so configurations stored by older versions of the framework
continue to load even after new fields are added. \texttt{config\_hash()}
serialises the configuration with \texttt{json.dumps(\dots, sort\_keys=True)}
and returns the first eight characters of the SHA-256 digest as a stable,
machine-independent fingerprint.

\subsection{Monte Carlo Engine Interface}
\label{sec:impl-engine-interface}

All engines are subclasses of the abstract base class \texttt{MonteCarloEngine}
defined in \texttt{benchmarking/engines/base.py}:

\begin{verbatim}
class MonteCarloEngine:
    name: str
    language: str
    backend: str

    def supports(self, workload_type: str) -> bool: ...
    def supported_ad_modes(self) -> tuple[str, ...]: ...
    def run(self, config: WorkloadConfig,
            ad_mode: str = "none") -> tuple[float, dict | None]: ...
\end{verbatim}

The class attributes \texttt{name}, \texttt{language} and \texttt{backend} are
the values written to the corresponding columns of \texttt{benchmarks.db},
enabling later queries such as ``all Rust runs'' or ``all GPU runs''. The
\texttt{run()} method is the only computational entry point: it accepts a
fully-validated \texttt{WorkloadConfig} and an AD mode string, and returns
$(\text{price},\ \text{greeks})$. The \texttt{greeks} component is either
\texttt{None} (for \texttt{ad\_mode = "none"}) or a dictionary mapping Greek
names (\texttt{"delta"}, \texttt{"vega"}, \texttt{"rho"}, \ldots) to floats.

\subsection{NumPy Engine}
\label{sec:impl-numpy}

\texttt{CPUMonteCarloEngine} (\texttt{benchmarking/engines/mc\_cpu.py}) is the
single-threaded NumPy reference implementation. For the European workload it
draws all $M$ standard normals in one call to
\texttt{np.random.default\_rng(seed).standard\_normal(M)}, evaluates the closed
form for $S_T$ vectorially and computes the discounted mean payoff. For the
local-volatility workload it runs the $N$-step log-Euler loop in Python, with
each iteration operating on a NumPy array of $M$ paths and using the stable
softplus

\begin{verbatim}
def softplus(z):
    return np.maximum(z, 0.0) + np.log1p(np.exp(-np.abs(z)))
\end{verbatim}

The same module exposes Black--Scholes closed-form helpers for price, Delta,
Vega and Rho, which the runner calls to populate the analytical reference
columns of every database row.

\subsection{JAX Engine}
\label{sec:impl-jax}

\texttt{JAXMonteCarloEngine} (\texttt{benchmarking/engines/mc\_jax.py}) is the
only engine that implements AD. The European kernel is declared at module level
to maximise compilation cache hits:

\begin{verbatim}
@functools.partial(jax.jit, static_argnames=("option_type",))
def _european_price(S0, r, sigma, K, T, Z, option_type):
    sqrtT = jnp.sqrt(T)
    ST = S0 * jnp.exp((r - 0.5 * sigma**2) * T + sigma * sqrtT * Z)
    payoff = jnp.where(option_type == "call",
                       jnp.maximum(ST - K, 0.0),
                       jnp.maximum(K - ST, 0.0))
    return jnp.exp(-r * T) * jnp.mean(payoff)
\end{verbatim}

Reverse-mode Greeks for the European workload are obtained from a single
backward pass:

\begin{verbatim}
grad_fn = jax.jit(jax.grad(_european_price, argnums=(0, 1, 2)),
                  static_argnames=("option_type",))
delta, rho, vega = grad_fn(S0, r, sigma, K, T, Z, option_type)
\end{verbatim}

Forward-mode uses three \texttt{jax.jvp} calls, one per input direction, each
probing the corresponding basis tangent vector. The local-volatility kernel
expresses the $N$-step loop as \texttt{jax.lax.scan}:

\begin{verbatim}
def _step(carry, z):
    S, t = carry
    x = jnp.log(S / S0)
    raw = a0 + a1 * x + a2 * x * x + b1 * t
    sigma = sigma_min + softplus(raw)
    S_next = S * jnp.exp((r - 0.5 * sigma**2) * dt
                         + sigma * jnp.sqrt(dt) * z)
    return (S_next, t + dt), None

(S_T, _), _ = jax.lax.scan(_step, (S0_arr, 0.0), Z_matrix)
\end{verbatim}

\texttt{jax.lax.scan} produces a single loop in the XLA graph; tracing time and
peak memory remain constant in $N$, in contrast with a Python \texttt{for} loop
which would unroll the graph to size $\Theta(N)$ and substantially increase
compile time and memory for $N=252$.

The AD overhead is computed by the runner as the ratio of the mean timed
runtime of the AD run to that of an identically-configured no-AD baseline run.
Crucially, both legs use the same warmup count, so XLA cache state is
equivalent and the ratio reflects only the differentiation cost.

\subsection{C++ Engine}
\label{sec:impl-cpp}

The native implementation lives in
\texttt{benchmarking/cpp/cpu\_engine.cpp}~/~\texttt{cpu\_engine.hpp} and is exposed
to Python via pybind11 bindings in
\texttt{benchmarking/cpp/bindings/pybind\_module.cpp}. The Python wrapper
\texttt{benchmarking/engines/mc\_cpp.py} loads the compiled module lazily and
falls back to a not-available stub if the import fails.

Parallelisation is via OpenMP. Each thread owns its own
\texttt{std::mt19937\_64} generator, seeded as \texttt{seed + thread\_id} to
guarantee statistical independence across threads while remaining reproducible
for any fixed thread count. The European inner loop is:

\begin{verbatim}
double sum = 0.0;
#pragma omp parallel reduction(+:sum)
{
    int tid = omp_get_thread_num();
    std::mt19937_64 rng(seed + tid);
    std::normal_distribution<double> N01(0.0, 1.0);
    #pragma omp for
    for (int i = 0; i < M; ++i) {
        double Z = N01(rng);
        double ST = S0 * std::exp((r - 0.5 * sigma * sigma) * T
                                  + sigma * std::sqrt(T) * Z);
        sum += std::max(option_type == 0 ? ST - K : K - ST, 0.0);
    }
}
double price = std::exp(-r * T) * sum / M;
\end{verbatim}

The local-volatility kernel reuses the same RNG strategy and implements the
log-Euler loop with the numerically stable softplus
\texttt{(raw > 0 ? raw : 0) + std::log1p(std::exp(-std::abs(raw)))}, identical
to the Python and Rust formulations.

\subsection{Rust Engine}
\label{sec:impl-rust}

The native implementation lives in \texttt{benchmarking/rust/src/lib.rs} and is
exposed to Python via PyO3 bindings, built with \texttt{maturin develop
--release}. The Python wrapper \texttt{benchmarking/engines/mc\_rust.py} mirrors
the C++ wrapper: lazy import, graceful fallback, identical engine interface.

The Rust implementation uses \texttt{SmallRng} (Xoshiro128++) from the
\texttt{rand} crate for high-throughput sampling, with each Rayon worker thread
seeded as \texttt{seed XOR thread\_index}. The simulation is expressed as a
data-parallel iterator over the path range:

\begin{verbatim}
let price: f64 = (0..M).into_par_iter()
    .map_init(
        || SmallRng::seed_from_u64(seed ^ rayon::current_thread_index()
                                            .unwrap_or(0) as u64),
        |rng, _| {
            let z: f64 = StandardNormal.sample(rng);
            let s_t = s0 * ((r - 0.5 * sigma * sigma) * t
                            + sigma * t.sqrt() * z).exp();
            (s_t - k).max(0.0)
        })
    .sum::<f64>() * (-r * t).exp() / (M as f64);
\end{verbatim}

The softplus is implemented inline as
\texttt{z.max(0.0) + (-z.abs()).exp().ln\_1p()}, matching the Python, C++ and JAX
formulations bit-for-bit.

\subsection{CUDA Engine}
\label{sec:impl-cuda}

\texttt{CUDAMonteCarloEngine} (\texttt{benchmarking/engines/mc\_cuda.py}) uses
PyCUDA to compile and launch CUDA kernels at runtime. At import time, the engine
runs a minimal 128-path trial; if CUDA initialisation, kernel compilation, or
the trial run fails, the engine sets its availability flag to \texttt{False} and
is excluded from \texttt{/api/engines}.

The European kernel uses one thread per Monte Carlo path:

\begin{verbatim}
__global__ void price_kernel(double S0, double r, double sigma,
                             double K, double T, double sqrtT,
                             unsigned long long seed,
                             int option_type, int M,
                             double *payoffs) {
    int idx = blockIdx.x * blockDim.x + threadIdx.x;
    if (idx >= M) return;
    curandState state;
    curand_init(seed, idx, 0, &state);
    double Z = curand_normal_double(&state);
    double ST = S0 * exp((r - 0.5 * sigma * sigma) * T + sigma * sqrtT * Z);
    double payoff = option_type == 0
                    ? fmax(ST - K, 0.0) : fmax(K - ST, 0.0);
    payoffs[idx] = exp(-r * T) * payoff;
}
\end{verbatim}

Using the thread index as the \texttt{curand\_init} sequence guarantees
statistical independence regardless of block or grid layout. The host averages
the resulting payoff array with \texttt{np.mean}.

Pathwise forward-mode Delta is implemented in a second kernel. Using the
identity $\partial S_T / \partial S_0 = S_T / S_0$ for GBM, the pathwise
Delta integrand is

\[
\mathbf{1}_{\{S_T > K\}} \cdot \frac{S_T}{S_0} \qquad (\text{call}), \qquad
-\mathbf{1}_{\{S_T < K\}} \cdot \frac{S_T}{S_0} \qquad (\text{put}),
\]

and the engine returns
$\Delta = e^{-rT}\,\mathbb{E}[d(\text{payoff})/dS_0]$. Note that this is a
\emph{pathwise} estimator, not a finite-difference approximation; it produces
unbiased Greeks with the same variance characteristics as a single-pass
forward-mode AD.

\subsection{Benchmark Runner}
\label{sec:impl-runner}

\texttt{BenchmarkRunner} (\texttt{benchmarking/runner.py}) is a small class.
Construction takes an engine instance, an optional display name and the warmup
and repetition counts. The \texttt{run()} method implements the protocol of
Section~\ref{sec:runner-layer}. Timed repetitions use \texttt{time.perf\_counter()}
which provides nanosecond resolution on Linux and macOS. The runner separately
computes the analytical reference values (by calling the Black--Scholes helpers
in \texttt{mc\_cpu.py}) and the absolute and relative errors for the price and
each Greek, so that the database row produced is self-contained.

When \texttt{ad\_mode != "none"} the runner first performs the AD timed run,
then a matched no-AD baseline run using the same engine, same configuration,
same warmup count and same repetition count. This ensures that the AD overhead
ratio reflects only the cost of differentiation, not the cost of XLA
recompilation or any other one-time setup.

\subsection{Metrics Collection}
\label{sec:impl-metrics}

The runner collects four classes of metrics:

\begin{itemize}
    \item \textbf{Timing:} mean, standard deviation, min, max of the timed
          repetitions, in milliseconds, plus the derived
          \texttt{throughput\_paths\_per\_sec = M / mean\_runtime\_s}.
    \item \textbf{AD overhead:} the matched baseline mean runtime
          (\texttt{baseline\_mean\_ms}) and the overhead ratio
          $\overline{t}_{\text{AD}}/\overline{t}_{\text{baseline}}$
          (\texttt{ad\_overhead\_ratio}).
    \item \textbf{Memory:} process RSS sampled before and after the timed
          repetitions via \texttt{psutil}; the recorded value is the difference.
    \item \textbf{Numerical accuracy:} absolute and relative error for the
          price and for each Greek against the Black--Scholes reference, when
          applicable.
\end{itemize}

Environment metadata is captured once per \texttt{run()} call: the wall-clock
timestamp, the Python version and implementation, the platform string, the CPU
model and architecture, the logical CPU count, total RAM in GB, the NumPy
version and the JAX version.

\subsection{Storage Layer}
\label{sec:impl-storage}

\texttt{BenchmarkDB} (\texttt{benchmarking/db.py}) is a thin SQLite wrapper. A
new connection is opened for each call with \texttt{check\_same\_thread=False},
and an explicit \texttt{threading.Lock} serialises write operations. The database
is created in WAL mode the first time it is opened. The schema is described in
detail in Section~\ref{sec:storage-layer}; it is defined declaratively in the
\texttt{\_FULL\_SCHEMA\_COLUMNS} constant, and any missing columns are added
automatically at startup via \texttt{PRAGMA table\_info} followed by
\texttt{ALTER TABLE}.

Two write paths coexist. Experiment scripts use \texttt{store\_run\_full()},
which writes a fully-populated \texttt{completed} row in a single transaction.
The REST API uses the three-step lifecycle
\texttt{create\_run()} $\to$ \texttt{mark\_running()} $\to$
\texttt{mark\_completed()} (or \texttt{mark\_failed()}) to support asynchronous
execution from a background worker thread.

\subsection{REST API}
\label{sec:impl-api}

The Flask application (\texttt{benchmarking/api/server.py}) registers the seven
endpoints listed in Table~\ref{tab:api-endpoints} and starts a single daemon
background thread on startup. The worker polls the \texttt{runs} table for
\texttt{pending} rows in FIFO order, dispatches each through the appropriate
engine via \texttt{BenchmarkRunner.run()}, and updates the row with the result.
Engine availability is probed once at startup: C++ and CUDA engines that fail
to load or fail their probe runs are removed from the registry.

The API returns JSON in all cases and uses standard HTTP status codes. Flask-CORS
is configured to allow requests from the Vite development server (port 5173) so
that the frontend can be developed locally without a reverse proxy.

\subsection{Frontend / Dashboard}
\label{sec:impl-frontend}

The frontend is a React~+~TypeScript single-page application built with Vite
and styled with Material UI. The directory layout follows the conventional
\texttt{src/pages/}, \texttt{src/components/} and \texttt{src/api/} structure.
A typed HTTP client in \texttt{src/api/client.ts} wraps every API endpoint with
a TypeScript function whose return type mirrors the corresponding Python
dataclass.

Pages:

\begin{itemize}
    \item \texttt{SimulatePage} --- submits a new run via POST
          \texttt{/api/runs} and polls GET \texttt{/api/runs/:id} until the
          status becomes \texttt{completed} or \texttt{failed}.
    \item \texttt{ComparePage} --- consumes \texttt{/api/compare\_matrix} to
          render a multi-engine, multi-AD-mode comparison grid.
    \item \texttt{HistoryPage} --- filterable table backed by GET
          \texttt{/api/runs}.
    \item \texttt{SummaryPage} --- consumes \texttt{/api/summary}.
    \item \texttt{AnalysisPage} --- focused view of AD overhead and Greek
          accuracy.
\end{itemize}

Components are factored to be workload-agnostic: \texttt{SimulationForm} is
generated dynamically from the workload \texttt{SCHEMA} returned by the API, so
adding a new workload class (or even a new field) requires \emph{no} frontend
changes.

\section{Supported Workloads}
\label{sec:impl-workloads}

Three workloads are implemented and validated end-to-end:

\begin{itemize}
    \item \texttt{european} (\texttt{EuropeanOptionConfig}): single-step exact
          GBM. Supported by all five engines. Used for AD overhead, Greek
          accuracy, language comparison and throughput scaling experiments.
    \item \texttt{european\_local\_vol} (\texttt{EuropeanLocalVolConfig}):
          $N$-step log-Euler simulation of GBM with a parametric
          local-volatility surface. Supported by NumPy, JAX, C++ and Rust;
          CUDA support is not implemented in the current iteration. Used to
          stress-test AD on a deeper computation graph and to expose the cost
          of multi-step over single-step simulation.
    \item \texttt{asian} (\texttt{AsianOptionConfig}): arithmetic-average
          path-dependent option workload. Supported by the Python and JAX
          engines in the final cloud-profiler experiments. Used to check that
          backend rankings and AD overheads are not inferred only from
          terminal-payoff European workloads.
\end{itemize}

\section{AD Integration}
\label{sec:impl-ad}

Automatic differentiation is implemented exclusively in the JAX and CUDA
engines. The JAX engine supports both forward and reverse mode for both
workloads; the CUDA engine implements only forward-mode pathwise Delta on the
European workload.

For the European workload the AD output is a three-tuple
$(\Delta, \rho, \mathrm{Vega}) = (\partial P/\partial S_0,\ \partial P/\partial r,\ \partial P/\partial \sigma)$.
Forward mode is implemented as three \texttt{jax.jvp} calls (one per input
direction), and reverse mode as a single \texttt{jax.grad} call with
\texttt{argnums=(0,1,2)}.

For the local-volatility workload the AD output is a six-tuple
$(\partial P/\partial S_0,\ \partial P/\partial a_0,\ \partial P/\partial a_1,\ \partial P/\partial a_2,\ \partial P/\partial b_1,\ \partial P/\partial \sigma_{\min})$.
Forward mode therefore requires six \texttt{jax.jvp} calls, while reverse mode
again computes the full gradient in a single backward pass.

Both modes feed into the runner's overhead-ratio computation. Because the
gradient and the baseline price share the same forward computation, the
overhead ratio is a clean measure of the marginal cost of differentiation.

\section{GPU Support}
\label{sec:impl-gpu}

CUDA support is implemented for the European workload via PyCUDA. The kernel
launch configuration is determined at runtime: the host queries the device for
the maximum number of threads per block and chooses the largest power of two
not exceeding it, then computes the number of blocks as
$\lceil M / \text{threads\_per\_block} \rceil$. The host allocates a single
device array of size $M$ for the payoffs and reduces it to a scalar with
\texttt{np.mean} after copying back. The CUDA engine is probed at server
startup with a 128-path trial run; if the probe fails (no GPU, no CUDA
toolkit, PyCUDA import failure, etc.) the engine is silently excluded.

JAX inherits GPU support transparently when JAX is installed with CUDA wheels.
The same JAX engine code paths execute on GPU without modification; the
benchmark database simply records a different platform string. Switching
between CPU and GPU JAX is therefore a deployment decision rather than a code
change.

\section{Challenges}
\label{sec:impl-challenges}

Several issues arose during implementation and shaped the final design.

\paragraph{XLA caching across runs.}
A naive initial JAX implementation defined the kernel inside the engine's
\texttt{run()} method. Each call created a fresh function object, defeating
the JIT cache: every timed repetition paid the full $\approx115$~ms compile
cost. Moving the kernel to module level reduced warm execution time to
$\approx0.6$~ms and made the JAX engine genuinely competitive on small $M$.

\paragraph{Triple AD passes.}
An earlier reverse-mode implementation invoked \texttt{jax.grad} three times
with \texttt{argnums=0,1,2}. The result was numerically correct but tripled
both the measured AD wall time and the AD overhead ratio. Consolidating into a
single \texttt{jax.grad(..., argnums=(0,1,2))} call eliminated this artefact;
this single change halved the reverse-mode overhead in early measurements and
brought the framework into line with how reverse-mode AD is deployed in
practice.

\paragraph{\texttt{for}-loop unrolling in JAX.}
The first local-volatility implementation used a Python \texttt{for} loop over
the $N=252$ Euler steps. Compilation time grew linearly in $N$ and peak memory
reached several gigabytes during tracing. Rewriting the loop as
\texttt{jax.lax.scan} produced a single fused loop in the XLA graph; compile
time and memory became effectively independent of $N$.

\paragraph{Numerically stable softplus.}
A naive softplus implementation $\log(1 + e^{z})$ overflows for $z \gtrsim 700$
in double precision and loses precision near $z = 0$. The stable form
$\max(z, 0) + \log(1 + e^{-|z|})$ is used identically in NumPy, JAX, C++ and
Rust, so any cross-engine numerical discrepancy can only be attributed to the
RNG or the floating-point order of operations, not to the activation function
itself.

\paragraph{Per-thread RNG reproducibility.}
Both the C++ and Rust engines use per-thread RNG streams seeded from the global
seed plus the thread index. Changing the thread count therefore changes the
RNG assignment and produces statistically equivalent but numerically different
results. This is recorded by storing \texttt{num\_threads} on every run, so
that runs at different thread counts are treated as distinct configurations.

\paragraph{Optional native components.}
Both the C++ and Rust engines require non-trivial native toolchains
(a C++17 compiler with OpenMP, or the Rust toolchain plus maturin) that are not
present in every environment. The engines are loaded lazily and degrade
gracefully: a missing build is reported only when the user explicitly requests
that engine via the API, and the corresponding rows simply do not appear in the
engine list.

\paragraph{Schema evolution.}
During development, several new columns were added to the database (notably
the AD-overhead, error-metric and cloud columns). Rather than write migration
scripts, the database opens with a self-migrating \texttt{\_init\_schema()}
that compares the existing columns against the authoritative list and
\texttt{ALTER TABLE}s any missing columns into existence. This has allowed every
historical database file to remain readable across schema changes.

\section{Cloud Metadata and Cost Support}
\label{sec:impl-cloud}

Cloud support is implemented as metadata attached to ordinary benchmark rows
rather than as a separate execution path. Experiment drivers can record the
cloud provider, instance type, region, zone, hourly price assumption, git
commit, and derived cost per run. This preserves the same framework invariants
as local execution: every cloud result still has a workload configuration, an
engine, an AD mode, timing statistics, numerical outputs, and environment
metadata.

The final cloud experiment uses Google Cloud instance metadata recorded in the
result files, including \texttt{n2-standard-4} in
\texttt{europe-west1-b}. Cost per run is derived from measured runtime and the
stored hourly price. The implementation deliberately avoids hard-coding a claim
that one instance is universally cheapest: prices and availability can change,
so the report treats cost as experiment metadata associated with a particular
run.

\section{Profiler Implementation}
\label{sec:impl-profiler}

The final profiler-versus-grid experiment is driven by
\texttt{experiments/run\_profiler\_vs\_grid.py}. The script constructs the
candidate grid, filters unsupported workload/engine/AD combinations, executes
probe runs, applies the old profiler and SHA selection policies, and records
which configurations would be promoted to full benchmarking. The report
artefact generator
\texttt{Final\_Year\_Project\_Report/evaluation/generate\_report\_assets.py}
then converts the stored results into the tables and figures used in
Chapter~\ref{chap:results}.

In the final \texttt{sha\_cloud\_profiler\_v4} experiment, the full-grid
baseline contains 16 candidate configurations evaluated at three full path
counts, for 48 full benchmark cells. The old probe-based profiler selects 30
full cells, while SHA selects 24. The SHA implementation records intermediate
promotion rounds, including the active configuration count at each probe path
budget. This metadata is used to generate the SHA progression plot and to
compute probe-budget savings.

The profiler implementation also computes decision-quality metrics against the
full grid: Pareto-frontier recovery, runtime regret, cost regret, and Spearman
rank correlation between probe and full-scale rankings. This is important
because the profiler is an optimisation of the benchmarking process, not a
replacement for evaluation. Its success is judged by whether it reduces
expensive full runs while preserving the decisions that the full grid would
have supported.

\section{Testing}
\label{sec:impl-testing}

The framework ships with $55$ tests across two pytest files in \texttt{tests/}.

\texttt{test\_system.py} covers:

\begin{itemize}
    \item Six parametric tests for configuration validation (invalid field
          values, rejected \texttt{option\_type}).
    \item Configuration serialisation: round-trip \texttt{to\_dict} /
          \texttt{from\_dict}, exclusion of \texttt{SCHEMA} from the serialised
          form, \texttt{config\_from\_dict} dispatch, and the expected
          \texttt{ValueError} on an unknown \texttt{workload\_type}.
    \item Engine correctness: NumPy and JAX pricing accuracy against
          Black--Scholes within tolerance.
    \item Runner statistics: structure of the timing record and
          reproducibility under fixed seed.
    \item Database round-trips: \texttt{create\_run},
          \texttt{mark\_running}, \texttt{mark\_completed}, \texttt{get\_run},
          \texttt{get\_all\_runs}, \texttt{summary}.
    \item API endpoint contracts: status codes and response structure for all
          seven endpoints.
    \item An end-to-end flow from configuration through engine and runner to
          database insert and API read-back.
\end{itemize}

\texttt{test\_ad\_framework.py} covers:

\begin{itemize}
    \item JAX engine executes without error for each of
          \(\texttt{ad\_mode} \in \{\texttt{none}, \texttt{forward}, \texttt{reverse}\}\)
          on the European
          workload.
    \item The standalone differentiated price function exposes the expected
          JAX function signature.
    \item Analytical Greek computations agree with their closed-form
          definitions.
    \item Computed gradients fall within the expected analytical bounds.
\end{itemize}

\section{Dependencies}
\label{sec:impl-deps}

Core Python dependencies: \texttt{numpy>=1.26}, \texttt{scipy>=1.11},
\texttt{jax[cpu]>=0.4.28}, \texttt{flask>=3.0}, \texttt{flask-cors>=4.0},
\texttt{pytest>=8.0}, \texttt{pybind11>=2.11} (for the C++ build). \texttt{psutil}
is optional and enables RSS-based memory profiling when present. The Rust
engine requires the Rust toolchain (\texttt{rustup}), \texttt{maturin>=1.0},
\texttt{PyO3}, and the \texttt{rand}, \texttt{rand\_distr} and \texttt{rayon}
crates. The CUDA engine requires the NVIDIA CUDA toolkit and \texttt{PyCUDA}.
The frontend depends on \texttt{react}, \texttt{react-dom}, \texttt{@mui/material}
and a Vite-based toolchain.

\section{Summary}
\label{sec:impl-summary}

The implementation realises the design of Chapter~\ref{chap:system-design} in
roughly $4{,}000$ lines of Python plus native C++ and Rust extensions and a
TypeScript frontend. Five engines (NumPy, JAX, C++, Rust, CUDA), two workloads
(European GBM, European local volatility) and three AD modes (none, forward,
reverse) are supported, with full SQLite persistence, a REST API, a React
dashboard and a $55$-test pytest suite. The next chapter describes the
experimental protocol applied to this implementation.
